\section*{A Life Made Legible by Work}
\addcontentsline{toc}{section}{A Life Made Legible by Work}

Thomas Aquinas left an immense intellectual record and almost no continuous
account of himself.  His questions, commentaries, sermons, disputed
determinations, and unfinished syntheses show a Dominican friar thinking in
public: receiving authorities, distinguishing meanings, meeting objections,
and trying to lead created intelligence toward God.  They do not disclose a
modern private life.  The scenes that have made him most familiar---the silent
schoolboy called an ox, the captive resisting temptation, the crucifix praising
his writing, and the scholar declaring his books straw---come principally from
canonization testimony and early hagiography assembled decades after his death.

This life therefore begins with a double fact.  Thomas is unusually accessible
as an author and unusually mediated as a personality.  Contemporary charters,
university and papal acts, the chronology of his works, and testimony from
people who knew him secure a remarkable itinerary through the schools and
convents of thirteenth-century Europe.  William of Tocco and the canonization
inquiries preserve human memories that no other archive supplies, but they also
arrange those memories to demonstrate sanctity.  Later Catholic reception
confesses Thomas as saint and Doctor of the Church; that judgment neither turns
every anecdote into eyewitness history nor makes every proposition timeless.

The governing question is how a younger son from a frontier aristocratic
household became a Friar Preacher, master of sacred doctrine, and author whose
account of creation, grace, reason, virtue, and beatitude could outlive the
institutional conflicts that formed it.  The answer is not the story of a lone
genius escaping his century.  It is the story of a vowed religious working
inside households, schools, courts, scriptoria, controversies, friendships,
and worship---and of a Church that received his work through argument as well
as praise.
